\documentclass[11pt,a4paper]{article}

\usepackage[margin=2.4cm]{geometry}
\usepackage{amsmath,amssymb}
\usepackage{booktabs}
\usepackage{graphicx}
\usepackage{microtype}
\usepackage{siunitx}
\usepackage{caption}
\usepackage{placeins}

\setlength{\parindent}{0pt}
\setlength{\parskip}{0.55em}
\captionsetup{font=small,labelfont=bf}
\sisetup{round-mode=places,round-precision=4}
\newcommand{\taskfigure}[2]{%
  \IfFileExists{#2}{%
    \includegraphics[width=#1\linewidth]{#2}%
  }{%
    \fbox{\parbox[c][4cm][c]{0.86\linewidth}{\centering
      Generated figure is available when this file is compiled from the
      \texttt{report} directory.}}%
  }%
}

\begin{document}

\section*{1.1 Part A: Whitening and autocovariance matrices}

\subsection*{1.1.1 Conventions and implementation}

The observations are stored as the rows of an $n\times p$ matrix $X$.  Each variable is
centred using its full-sample column mean,
\[
  X_c=X-\mathbf{1}_n\bar{x}^{\mathsf T},
  \qquad
  S_0=\frac{1}{n-1}X_c^{\mathsf T}X_c.
\]
If $S_0=VDV^{\mathsf T}$, the symmetric whitening matrix and whitened data are
\[
  W=S_0^{-1/2}=VD^{-1/2}V^{\mathsf T},
  \qquad Y=X_cW.
\]
Eigenvalues below $10^{-10}$ times the largest eigenvalue are treated as numerically zero,
because their inverse square roots would be unstable.  No additional marginal scaling is
applied: scaling is already included in $W$.

For lag $h\geq1$, the implementation uses
\[
  \widehat C_h
  =\frac{1}{n-h}\sum_{t=1}^{n-h}y_ty_{t+h}^{\mathsf T}.
\]
The mean is not re-estimated inside each lag window.  The default matrix used for subsequent
orthogonal diagonalisation is the symmetric version
$\widehat C_h^{(s)}=(\widehat C_h+\widehat C_h^{\mathsf T})/2$.  Thus lagging pairs row $t$
with row $t+h$, and normalization uses the number of available pairs, $n-h$.

\subsection*{1.1.2 Purpose of whitening}

Whitening removes contemporaneous scale and cross-covariance, so the transformed sample has
lag-zero covariance $I_p$.  In the second-order source-separation model this changes the
unknown mixing matrix into an orthogonal one.  The remaining directions can therefore be
estimated by orthogonal diagonalisation of the lagged autocovariance matrices.  Whitening
does \emph{not} turn the series into white noise: it removes dependence between variables at
lag zero, while serial dependence at non-zero lags remains in $\widehat C_h$ and supplies the
signal information.

\subsection*{1.1.3 Numerical verification}

For the supplied data, $n=1000$ and $p=10$.  The numerical check gave
\[
 \max_j|\bar y_j|=6.49\times10^{-14},\qquad
 \max_{j,k}|\{\widehat{\operatorname{Cov}}(Y)-I_{10}\}_{jk}|
 =3.55\times10^{-15}.
\]
The whitened variances ranged from $0.999999999999997$ to $1$, and the largest absolute
off-diagonal covariance was $2.58\times10^{-15}$.  These values confirm the expected effect
to numerical precision.  The non-zero Frobenius norms of the lagged matrices (for example,
$0.3653$ at lag 1 and $0.5575$ at lag 2) also show that whitening does not force
non-zero-lag covariance estimates to zero.

\begin{figure}[htbp]
  \centering
  \taskfigure{0.88}{../output/figures/task1_whitening_covariance.png}
  \caption{Original contemporaneous correlation and covariance after whitening.}
\end{figure}

\FloatBarrier
\section*{1.2 Part B: Bootstrap for dependent observations}

\subsection*{1.2.1 Fixed-block bootstrap}

For block length $L$, a block start is sampled uniformly with replacement from
$\{1,\ldots,n-L+1\}$.  The corresponding $L$ consecutive multivariate rows are retained in
their original order.  This is repeated $\lceil n/L\rceil$ times, the blocks are concatenated,
and the result is truncated to $n$ rows.  The tuning parameter $L$ controls the amount of
local dependence retained within a block.

\subsection*{1.2.2 Stationary bootstrap}

The first index is uniform on $\{1,\ldots,n\}$.  At each later time, a new uniform index is
selected with probability $r=1/L$; otherwise the preceding index is advanced by one, with
circular wrap-around from $n$ to 1.  Hence the block length $K$ is geometric,
\[
  \Pr(K=k)=r(1-r)^{k-1},\qquad E(K)=\frac{1}{r}=L.
\]
Unlike the fixed-block method, block lengths are random and circular continuation avoids a
special boundary.  In 100 generated samples, target expected lengths 10, 20 and 40 produced
mean observed lengths 9.94, 20.20 and 40.94, respectively.

\subsection*{1.2.3 Why iid resampling is unsuitable}

Sampling individual rows independently destroys their temporal ordering and therefore makes
non-zero-lag autocovariances approximately zero.  The resulting bootstrap sample no longer
represents the statistic's variability under the original dependent process.  Block methods
resample consecutive multivariate observations and retain short-range serial and cross-series
dependence within blocks.

\subsection*{1.2.4 Effect of block length}

Dependence preservation was measured by the mean absolute difference between the original
and bootstrap sample autocorrelations, averaged over all 10 variables and lags 1--20.  Each
entry below summarizes 100 bootstrap samples.

\begin{table}[htbp]
  \centering
  \caption{Mean absolute ACF error; smaller values indicate better local-dependence preservation.}
  \begin{tabular}{lrr}
    \toprule
    Method & Mean error & Standard deviation \\
    \midrule
    iid & 0.0460 & 0.0017 \\
    Fixed, $L=10$ & 0.0381 & 0.0021 \\
    Fixed, $L=20$ & 0.0320 & 0.0019 \\
    Fixed, $L=40$ & 0.0270 & 0.0025 \\
    Stationary, $E(K)=10$ & 0.0336 & 0.0020 \\
    Stationary, $E(K)=20$ & 0.0295 & 0.0024 \\
    Stationary, $E(K)=40$ & 0.0269 & 0.0027 \\
    \bottomrule
  \end{tabular}
\end{table}

The iid method has the largest error.  Longer blocks preserve the observed ACF more closely,
although they also yield fewer effectively independent blocks and slightly greater Monte
Carlo variability.  The expected block length is therefore a bias--variance tuning choice,
not a parameter that should automatically be made as large as possible.

\begin{figure}[htbp]
  \centering
  \taskfigure{0.78}{../output/figures/task1_bootstrap_acf_comparison.png}
  \caption{Example ACF comparison for the variable with the largest absolute lag-1 sample ACF.}
\end{figure}

\begin{figure}[htbp]
  \centering
  \taskfigure{0.9}{../output/figures/task1_bootstrap_acf_error.png}
  \caption{ACF error over 100 bootstrap samples for each method and block length.}
\end{figure}

\FloatBarrier
\section*{1.3 Part C: The ladle quantities}

\subsection*{1.3.1 Eigenvalue-based quantity $\phi$}

For the lag set $\mathcal T$, an orthogonal matrix $U$ approximately jointly diagonalises the
whitened autocovariance matrices.  The component strengths are
\[
  s_j=\sum_{h\in\mathcal T}
  \left[\operatorname{diag}\!\left(U^{\mathsf T}\widehat C_h^{(s)}U\right)_j\right]^2,
\]
ordered so that $s_0\geq\cdots\geq s_{p-1}$.  The normalized scree contribution is
\[
  \phi(k)=\frac{s_k}{1+\sum_{\ell=0}^{p-1}s_\ell},
  \qquad k=0,\ldots,p-1.
\]
Large values represent directions with serial dependence; noise directions should have
strengths close to zero.  With several lags, $U$ is obtained using \texttt{JADE::rjd}; the
strength calculation and ordering are implemented directly.

\subsection*{1.3.2 Bootstrap-variability quantity $f$}

Let $\widehat B_k$ and $B^*_{k,b}$ contain the first $k$ ordered directions from the original
and $b$th bootstrap samples.  The subspace instability is
\[
  f_0(k)=\frac{1}{B}\sum_{b=1}^{B}
  \left\{1-\left|\det\!\left(\widehat B_k^{\mathsf T}B^*_{k,b}\right)\right|\right\},
  \quad k=1,\ldots,p-1,
  \qquad f_0(0)=0,
\]
and
\[
  f(k)=\frac{f_0(k)}{1+\sum_{\ell=0}^{p-1}f_0(\ell)}.
\]
The determinant measures agreement of the two $k$-dimensional subspaces and is unaffected by
sign changes of individual directions.  Stable signal subspaces have small $f(k)$; instability
increases after noise directions enter the subspace.

\subsection*{1.3.3 Combination and dimension estimate}

The ladle criterion and estimator are
\[
  g(k)=f(k)+\phi(k),
  \qquad
  \widehat q=\underset{k\in\{0,\ldots,p-1\}}{\arg\min}\;g(k).
\]
Thus $\phi$ decreases when weak or noise directions are included, while $f$ increases when the
estimated subspace becomes unstable.  Their sum balances these two pieces of information.

\subsection*{1.3.4 Simulation check}

The code was checked using Setting 1 of Nordhausen and Virta (2018): five Gaussian AR/MA
signal series and five Gaussian white-noise series, so $p=10$ and the true dimension is
$q=5$.  The check used $n=1000$, lags $1{:}12$, 100 bootstrap replications, and block length
20.  Both bootstrap versions returned $\widehat q=5$.  The largest whitening identity error
was $3.22\times10^{-15}$, the largest orthogonality error in $U^{\mathsf T}U$ was
$2.22\times10^{-15}$, and the joint off-diagonal squared-energy ratio was $0.0718$.  These
checks establish correct whitening, orthogonal directions, useful approximate joint
diagonalisation, and recovery of the known signal dimension.

\begin{figure}[htbp]
  \centering
  \taskfigure{0.92}{../output/figures/task1_ladle_criterion.png}
  \caption{Implemented $f$, $\phi$ and $g$ on simulated Setting 1 data.  The minimum is at the true dimension $q=5$ for both bootstrap methods.}
\end{figure}

\end{document}
